\section{Trabalhos Relacionados}

\subsection{Sistemas multiagente com LLMs em finanças}

Trabalhos recentes vêm explorando agentes baseados em LLMs para análise
financeira e para tomada de decisão em trading. TradingAgents organiza
diferentes papéis (analistas de notícias, analistas técnicos, gerentes de
risco e traders) em fluxos de comunicação que terminam em decisões
operacionais \citep{tradingagents2024}. FinCon propõe um sistema com
reforço verbal conceitual para apoiar decisões financeiras complexas,
explorando como agentes podem revisar suas próprias conclusões a partir de
sinais contextuais \citep{fincon2024}. AlphaAgents foca especificamente em
construção de carteiras de ações usando múltiplos agentes baseados em LLMs
\citep{alphaagents2025}, e HedgeAgents discute uma abordagem multiagente
voltada ao suporte à decisão em fundos \citep{hedgeagents2025}. Em comum,
essas propostas decompõem a tarefa em papéis especializados e combinam os
resultados desses papéis em uma decisão final.

O foco predominante, contudo, está em seleção de ativos ou em decisão de
compra e venda. O presente trabalho se posiciona em outro recorte: tratar o
problema como diagnóstico personalizado de uma carteira já existente,
condicionado ao perfil declarado pelo usuário. A unidade de avaliação
deixa de ser apenas a plausibilidade de uma tese de investimento e passa
a incluir a adequação da classificação aos diferentes perfis e a
capacidade do sistema de indicar incerteza quando faltam dados. Isso
afeta diretamente as métricas escolhidas: além de variáveis operacionais,
importam sensibilidade à personalização e cobertura dos fatores de risco.

\subsection{Debate multiagente e baselines fortes}

Debate entre múltiplos agentes baseados em LLMs tem sido proposto como
estratégia para melhorar raciocínio, robustez e verificação cruzada de
respostas. Evidências recentes, porém, mostram que parte dos ganhos
atribuídos ao debate desaparece quando se utilizam baselines de agente
único com orçamento computacional equivalente, com prompts melhores e
self-consistency \citep{zhang2025stop}. O ponto é especialmente relevante
em domínios com alto custo cognitivo, como o financeiro, onde acrescentar
chamadas a uma LLM se traduz em mais latência e mais custo por análise.

Por isso, este trabalho adota baselines fortes de agente único como
referência principal de comparação. Em particular, considera-se uma LLM
única com raciocínio estruturado e self-consistency, com orçamento
ajustado ao número de chamadas dos demais baselines, de modo a evitar que
diferenças observadas decorram apenas da maior alocação de computação em
sistemas multiagente. O sistema sem debate, com agentes especialistas e
moderador, entra como ponto intermediário e permite separar o efeito do
debate do efeito da decomposição em papéis.

\subsection{Personalização em diagnóstico financeiro}

Personalização aparece frequentemente como característica desejável em
aplicações financeiras com LLMs, mas raramente como dimensão de avaliação
em si. Em muitos sistemas, o perfil do usuário é passado como entrada sem
que se verifique, de forma controlada, se a saída de fato muda com esse
perfil. O risco é claro: respostas pouco condicionadas ao usuário, em que
o sistema repete diagnósticos genéricos sobre a carteira independente do
perfil declarado.

Aqui, a personalização é avaliada por meio de casos contrafactuais. A
mesma carteira é diagnosticada sob perfis conservador, moderado e
agressivo, com todas as demais entradas constantes. A expectativa é que a
classificação e a justificativa se desloquem na direção compatível com
cada perfil: uma carteira de alta volatilidade deve ser mais
frequentemente classificada como \textit{atenção} ou \textit{risco
elevado} sob perfil conservador do que sob perfil agressivo. A frequência
com que essa relação se sustenta passa a ser, então, uma métrica direta
de sensibilidade à personalização, e separa comportamento personalizado
de mera variação textual.
